\clearpage
\fancyhead[L]{\small COMET.}
\hypertarget{native-474-lower}{}
\begin{center}
\rule{0.55\linewidth}{0.4pt}\\[2pt]
Astronomische Nachrichten. Band XXII. Nr. 521. Seite 277. 278. 1844 November 28.\\[2pt]
\rule{0.55\linewidth}{0.4pt}
\end{center}
\begin{flushright}Göttingen 15. October und 11. November 1844.\end{flushright}

Von dem \emph{Cometen} habe ich am 6. October ein paar Vergleichungen mit Ceti 200 \textsc{Bode}, aus denen nur die Rectascension sich mit Sicherheit ableiten lässt.

\begin{center}
1844 October 6.\quad $10^{\rm u}\ 40^{\rm m}\ 32^{\rm s}$ m. Z.\quad Gerade Aufsteig. $18^\circ\ 50'\ 51''1$
\end{center}

Indem ich für die Sterne, mit welchen ich am 14. October den Cometen verglichen habe, in Folge der verschiedenen neueren Meridianbeobachtungen folgende scheinbare Positionen am Beobachtungstage zum Grunde lege:

\begin{center}
\setlength{\tabcolsep}{3pt}
\begin{tabular}{llll}
Stern 1. & \shortstack[l]{Gerade Aufsteigung\\$19^\circ\ 46'\ 4''0$} & \shortstack[l]{Abweichung\\$0^\circ\ 56'\ 51''4$} & südlich \\
\hfill 2. & \hfill $19\ 52\ 18.8$ & \hfill $0\ 44\ 45.5$ & \\
\end{tabular}
\end{center}
\noindent finde ich den Cometenort

\begin{center}
\setlength{\tabcolsep}{3pt}
\begin{tabular}{lll}
1844 October 14. & $9^{\rm u}\ 27^{\rm m}\ 6^{\rm s}$ m. Z. & Gerade Aufsteig. $20^\circ\ 16'\ 11''3$ \\
 & $9\ 32\ 7$ & Südl. Abweich. $0\ 44\ 30.3$ \\
\end{tabular}
\end{center}

\Needspace{7\baselineskip}
\SourceLeaf{475}
Dr. \textsc{Goldschmidt} hat für den Cometen folgende elliptische Elemente berechnet:

\begin{center}
Durchgang durch das Perihel 1844 Sept. 2.49152 mittl. Berliner Zeit.
\end{center}
{\small
\[
\left.\begin{array}{llr}
\text{Länge des Perihels} & \text{.\ .\ .\ .\ .} & 342^\circ\ 29'\ 44''9 \\
\text{Aufsteigender Knoten} & \text{.\ .\ .\ .\ .} & 63\ 48\ 55.2
\end{array}\right\}
\begin{array}{l}
\text{vom mittl. Aequin. 1844} \\
\text{Sept. 21.5 gezählt}
\end{array}
\]
}
\begin{center}
\setlength{\tabcolsep}{4pt}
\begin{tabular}{llr}
Neigung der Bahn & \mbox{.\ .\ .\ .\ .} & $2\ 55\ 1.9$ \\
\multicolumn{2}{l}{Logarithm der halben grossen Axe} & $0.4929151$ \\
Excentricität & \mbox{.\ .\ .\ .\ .\ .\ .\ .} & $0.6186103$ \\
\multicolumn{3}{l}{Bewegung rechtläufig.}
\end{tabular}
\end{center}

Die aus diesen Elementen folgende siderische Umlaufszeit ist 2400.339 Tage. Die Elemente stellen die 3 zum Grunde gelegten Beobb. (Aug. 23, Sept. 21, Oct. 15) genau dar; Dr. \textsc{Goldschmidt} wird demnächst sämmtliche bekannt gewordene Beobachtungen damit vergleichen; vorläufig hat er die Vergleichung mit den drei Beobb. vom 14. Oct., 1. Nov., 2. Nov. durchgeführt, wobei sich die Unterschiede ergeben haben:

\begin{center}
\begin{tabular}{lrrl}
 & in Ger. Aufst. & in Abweichung & \\
Oct. 14 & $-0''9$ & $-8''3$ & Gauss \\
Nov. 1 & $+8.7$ & $-10.3$ & Petersen \\
\hfill 2 & $-12.3$ & $-8.4$ & Rümker \\
\end{tabular}
\end{center}

Die folgende Ephemeride ist nach diesen Elementen berechnet.

\clearpage
\hypertarget{native-475-lower}{}
\begin{center}
\rule{0.55\linewidth}{0.4pt}\\[2pt]
Astronomische Nachrichten Band XXII. Nr. 526. Seite 353. 354. 1845 Januar 25.\\[2pt]
\rule{0.55\linewidth}{0.4pt}
\end{center}

Am 11. verglich Herr \textsc{Petersen} den Cometen mit zwei Sternen der Hist. Cél. p. 237 u. 295, deren scheinbare Oerter für die Beobachtungszeit er annahm

\begin{center}
\begin{tabular}{rr}
$19^{\rm u}\ 21^{\rm m}\ 2^{\rm s}85$ & $+42^\circ\ 1'\ 5''1$ \\
$22\ 20.95$ & $41\ 55\ 11.8$ \\
\end{tabular}
\end{center}
\noindent Er fand:

\begin{center}
Jan. 11.\quad $9^{\rm u}\ 16^{\rm m}\ 7^{\rm s}$ Alt. m. Z.\quad $289^\circ\ 37'\ 35''6$\quad $+41^\circ\ 58'\ 40''7$.
\end{center}

Aus zwei Vergleichungen mit denselben Sternen fand Herr Hofrath \textsc{Gauss}

\begin{center}
\begin{tabular}{lll}
Jan. 11. & $5^{\rm u}\ 35^{\rm m}\ 19^{\rm s}2$ Gött. m. Z. & $+41^\circ\ 54'\ 40''2$ \\
 & $5\ 39\ 29.0$ & AR. $289^\circ\ 41'\ 47''6$ \\
\end{tabular}
\end{center}

Die Declination beruht blos auf dem zweiten Sterne, die AR. auf dem Mittel aus beiden. Herr Hofrath \textsc{Gauss} wünscht, dass die Sterne in der unteren Culmination jetzt bestimmt werden mögen.

Am Schlusse dieses Artikels erhalte ich noch von Herrn Hofrath \textsc{Gauss} eine Göttinger Beobachtung, von Herrn \textsc{Rümker} folgende Elemente, die er früher geliefert haben würde, wenn nicht die Geschäfte der Navigationsschule, bei denen ihm mehr Hülfe zu wünschen wäre, jetzt fast alle seine Zeit in Anspruch nähmen, und von Herrn \textsc{Sievers} die Verbesserung seiner früheren Bahn.

\begin{center}
\setlength{\tabcolsep}{3pt}
\begin{tabular}{lll}
1845 Jan. 12. & $5^{\rm u}\ 51^{\rm m}\ 11^{\rm s}0$ M. Z. Göttingen & AR. $289^\circ\ 14'\ 45''5$ \\
 & $5\ 54\ 59.5$ & D. $42\ 22\ 15.4$ N. \\
\end{tabular}
\end{center}

\Needspace{7\baselineskip}
\SourceLeaf{476}
Zur Vergleichung dienten 3 Sterne, deren aus der Hist. Cél. abgeleitete Positionen, für den Beobachtungstag folgende sind:

\begin{center}
\begin{tabular}{rr}
$288^\circ\ 54'\ 5''4$ & $+42^\circ\ 15'\ 49''7$ \\
$289\ 1\ 5.6$ & $42\ 23\ 59.7$ \\
$290\ 7\ 35.2$ & $42\ 29\ 48.2$ \\
\end{tabular}
\end{center}

Der letzte dieser Sterne ist in \textsc{Hardings} Charte bedeutend zu nördlich eingezeichnet.
\begin{flushright}\textsc{Schumacher}.\end{flushright}
